\documentclass{article}
\usepackage[utf8]{inputenc}
\usepackage[T1]{fontenc}
\usepackage[a4paper]{geometry}		% Reduire les marges
\geometry{hmargin=3cm,vmargin=3.5cm}
\usepackage[english]{babel} 
\usepackage{amsfonts}
\usepackage{amsthm}
\usepackage{amsmath}
\usepackage{amssymb}
\usepackage{hyperref}
\usepackage{array}
\usepackage[svgnames]{xcolor}
\usepackage{xpatch}
\usepackage{tikz}
\usepackage{mathrsfs}
%\usepackage{graphicx}
%\usepackage{caption}
\usepackage[svgnames]{xcolor}
\usepackage{eurosym}
\usetikzlibrary{positioning,calc}
\renewcommand{\thesection}{\Roman{section}}
\xpatchcmd{\proof}{\itshape}{\normalfont\proofnameformat}{}{}
\newcommand{\proofnameformat}{\bfseries}
\theoremstyle{definition}
\newtheorem{defi}{Definition}[section]
\newtheorem{theo}{Theorem}[section]
\newtheorem{prop}[theo]{Proposition}
\newtheorem{lem}[theo]{Lemma}
\newtheorem{coro}[theo]{Corollary}
\newtheorem{exe}{Example}[section]
\newtheorem{rem}{Remark}[section]
\newtheorem{nota}{Notation}[section]
\renewcommand\qedsymbol{$\blacksquare$}
\numberwithin{equation}{section}

\DeclareMathOperator{\R}{\mathbb{R}}
\DeclareMathOperator{\Z}{\mathbb{Z}}
\DeclareMathOperator{\Q}{\mathbb{Q}}
\DeclareMathOperator{\N}{\mathbb{N}}
\DeclareMathOperator{\E}{\mathbb{E}}
\DeclareMathOperator{\C}{\mathcal{C}}

\title{Exo 7, corrigé}       
\author{Guillaume Woessner}
\date{}                       % sans date : celle du jour de compilation

%\pagestyle{headings}          % Pour mettre des entetes avec les titres
                              % des sections en haut de page
\sloppy                       % Ne pas faire deborder les lignes dans la marge

\begin{document}

\maketitle                    % Faire un titre utilisant les donnees
                              % passees : \title, \author et \date

\begin{abstract}
\begin{center}
Les théorèmes de convergence.
\end{center}
\end{abstract}

%\tableofcontents              % Table des matieres


Dans tout ce document, l'espace de probabilité filtré considéré est $(\Omega,\mathcal{F},(\mathcal{F}_n)_n,\mathbb{P})$, et on supposera que $(X_n)_n$ est une martingale adaptée à la filtration.


\section{Récap}

\paragraph{Exercice 3} Trouver un exemple où $X_n\mapsto X$ presque sûrement, mais où $\E[X_n]\nrightarrow \E[X]$.

\begin{proof}
Soit $(S_n)_n$ une marche aléatoire simple sur $\Z$, $X$ une variable aléatoire entière non-intégrable, et $\tau:=\inf\{n\geq 0 ~:~S_n=X\}$. Poser $X_n := S_{\tau\wedge n}$.
\end{proof}


\paragraph{Exercice 4} 
Supposons que $X_n\geq 0$ et que $\tau$ est un temps d'arrêt (potentiellement infini).\\
Comment peut-on définir $X_\tau$ ? Et si $(X_n)_n$ n'est pas forcément positive ?

\begin{proof}
Si $X_n\geq 0$, sur $\{\tau=\infty\}$ on peut poser $X_\tau=X_\infty=\lim_n X_n$, car la $\lim$ existe toujours d'après un corollaire du cours. Sans la condition de positivité on ne peut rien dire.
\end{proof}


\paragraph{Exercice 5} Démontrez le lemme sur le critère d'uniforme intégrabilité suivant :\\
Supposons que $X_n\mapsto X$ en $\mathbb{P}$ et que $X$ soit intégrable.\\
Montrez que la convergence a également lieu dans $L^1$ si et seulement si $(X_n)_n$ est U.I..\\
Remarquez que ceci correspond à une réciproque partielle de la propriété : convergence $L^1$ implique convergence en $\mathbb{P}$. 

\begin{proof}
Supposons d'abord que $(X_n)_n$ est U.I. et posons $\epsilon>0$ ainsi que le $K_\epsilon$ correspondant. On pourra supposer que $X=0$ (pourquoi?). On a alors
\begin{align*}
E|X_n| &\leq \E|X_n|\mathbf{1}_{|X_n|\leq \epsilon} + \E|X_n|\mathbf{1}_{\epsilon < |X_n|\leq K_\epsilon} + \E|X_n|\mathbf{1}_{K_\epsilon<|X_n|}\\
&\leq \epsilon + K_\epsilon \mathbb{P}(|X_n|>\epsilon) + \epsilon
\end{align*}
Passer à la $\limsup_n$ permet de conclure.\\
Supposons maintenant que la convergence a lieu dans $L^1$. On pourra supposer que $X=0$ (pourquoi?). Soit $\epsilon>0$. On sait que il existe $N_\epsilon$ tel que pour tout $n\geq N_\epsilon$ on ait $\E|X_n|\mathbf{1}_{|X_n|>a}\leq\E|X_n|\leq \epsilon$ pour tout $a>0$. Pour les $N_\epsilon-1$ premiers $X_n$ on peut choisir $K_\epsilon$ tel que pour tout $n\leq N_\epsilon-1$ on ait $\E|X_n|\mathbf{1}_{|X_n|>K_\epsilon}\leq \epsilon$. Poser $a=K_\epsilon$ dans la première partie permet de conclure.
\end{proof}


\paragraph{Exercice 6}
Soient $(X_n)_n$ et $(Y_n)_n$ deux martingales de carré intégrable.\\
Montrez que 
$$ \E X_nY_n - \E X_0Y_0 = \sum_{m=1}^n \E(X_m-X_{m-1})(Y_m-Y_{m-1}). $$
Remarquez que ça peut permettre de calculer la covariance de $X_n$ et $Y_n$.

\begin{proof}
Utilisez que 
$$\E[X_nY_n] = \dfrac{1}{2}\left(\E[(X_n+Y_n)^2]-\E[X_n^2]-\E[Y_n]^2\right).$$
\end{proof}

\section{Pour ceux qui ont vu les chaînes de Markov}

\paragraph{Exercice 7} 
Supposons que $(M_n)_n$ soit une chaîne de Markov homogène sur un espace dénombrable $S$ et dont le noyau de transition est noté $p~:~S\times S\to [0,1]$. On dit que $h~:~S\to \R$ est $p$-harmonique si pour tout $x\in S$ on a $h(x)=\sum_S p(x,y)h(y)$.\\
Montrer qu'une telle fonction $h$ bornée est $p$-harmonique si et seulement si $h(M_n)$ est une martingale adaptée à la filtration canonique de $(M_n)_n$.

\begin{proof}
Supposons d'abord que $(h(M_n))_n$ soit une martingale. On a alors par hypothèse
\begin{align*}
h(x)=\E[h(M_{n+1})~|~M_n=x]=\E_x[h(M_1)]=\sum_{y\in S}p(x,y)h(y).
\end{align*}
Ainsi $h$ est $p$-harmonique.\\
Le même type de raisonnement permet de conclure pour l'autre sens.
\end{proof}



\paragraph{Exercice 8} 
Supposons que $(M_n)_n$ soit une chaîne de Markov irréductible sur un espace dénombrable $S$. Soit $f~:~S\to\R^+$ une fonction non constante telle que pour tout $x\in S$,
$$\E_x[f(M_1)]\leq f(x).$$
Montrez que $M_n$ est transiente.

\begin{proof}
La condition implique que $(f(M_n))_n$ est une sur-martingale, positive par hypothèse. Ainsi on peut utiliser utiliser le théorème de convergence pour les sur-martingales positives : il existe $Y$ telle que $f(X_n)\mapsto Y$ presque sûrement.\\
On ne peut donc pas avoir que la chaîne de Markov est récurrente, sinon $X_n$ prendrait successivement toutes ses valeurs, et $f(X_n)$ ne pourrait pas converger.
\end{proof}




\end{document}












